\documentclass[10pt]{article}
\usepackage[margin=20mm]{geometry}
\usepackage{amsmath,amssymb,mathtools,fontspec,unicode-math,xurl,hyperref}
\IfFontExistsTF{DejaVu Serif}{\setmainfont{DejaVu Serif}}{\setmainfont{Latin Modern Roman}}
\setmathfont{latinmodern-math.otf}
\setlength{\parindent}{0pt}\setlength{\parskip}{6pt}\sloppy
\title{Formal initial review: M25-0076
The Assouad dimension of Kakeya sets in \(\mathbb R^3\)}
\author{Scientific review performed in this conversation}
\date{26 October 2026}
\begin{document}
\maketitle
\section*{Scope and initial outcome}
All 48 pages of the bound author manuscript were read, including references. Page numbers below are both the printed and one-based PDF numbers. The acquired `main.tex` was parsed completely for structure and dependencies and read directly at the principal disputed passages. The report applies to this v2 author manuscript. The journal reference attached to its metadata does not establish identity with the publisher’s version of record.\par
Initial outcome: Major. Two finding groups are established. First, the unrestricted twisted-projection argument in Section 7 exceeds the domains on which its geometric reduction and cinematic-function input are valid. There is an explicit horizontal-tube counterexample to the unrestricted Proposition 7.1 and Lemma 7.12, and an independent polynomial check of the unnormalized cinematic-family claim. Second, the parameter choice in the proof of Theorem 1.4 does not force the claimed contradiction. A complete local correction of that deduction is given below, conditional only on the stated Theorem 1.5.\par
ID | Confirmed claim or inference | Assessment\par
F1 | Proposition 7.1, p. 39; Lemma 7.12, pp. 42–44: direction-domain and cinematic-normalization restrictions are absent from the stated/generalized argument. | Major\par
F2 | Theorem 1.4 proof, p. 46: the measure-error exponent is too large for the stated dimension gap; the final exponent comparison is invalid. | Major\par
These findings do not constitute a counterexample to the main Kakeya dimension conclusions. The horizontal construction fails the stronger Convex Wolff Axioms used for the main theorem. It is excluded by an explicit transversality convention. The measure argument has a direct repair. What remains is to state the correct restricted projection input and verify that the main proof supplies its hypotheses, with the necessary quantitative losses.\par
The proof chain was traced through the two extremal parameters, the four-way structural alternative, plank enlargement, the generalized sticky Kakeya argument, twisted projections, and the final dimension deductions. Fifteen reviewer-written checks support specific calculations. They do not replace an all-scale proof or certify every imported result.\par
\section*{F1(a). Horizontal tubes refute the unrestricted projection claim}
Locations. Proposition 7.1, p. 39; Lemma 7.12 and (7.2), p. 42; the refinement-and-rotation reduction to lines \((a+ct,b+dt,t)\), p. 43. The tube definition on pp. 3 and 6 allows arbitrary directions. The restatement on p. 33 specifies a thickened line in a cube and does not explicitly impose a lower bound for its \(z\) component.\par
Construction. Fix a sufficiently large absolute integer \(M\). For arbitrarily small \(\delta\), take \(a,b\) in the \(M\delta\) lattice inside \([-1/8,1/8]\), and set \[v_b=\left(\frac{1-b^2}{1+b^2},\frac{2b}{1+b^2},0\right),\qquad
 L_{a,b}=\{(0,a,0)+t v_b:-1/2\leq t\leq1/2\}.\] Let \(T_{a,b}=N_\delta(L_{a,b})\) and \(Y(T)=T\). The vector \(v_b\) has norm one. Thus these are genuine unit-segment tubes. Their centers have norm at most \(1/8\), so the segments lie in \(B(0,5/8)\) and the tubes lie in \(B(0,1)\) for small \(\delta\). There are \(N\asymp\delta^{-2}\) tubes. Taking \(M\) large ensures essential distinctness, although Proposition 7.1 does not separately require it.\par
All-scale Tube Wolff check. For \(\delta\leq r\ll1\), if an \(r\)-tube contains one of these unit segments, its axis must be within \(O(r)\) in angle of the segment direction. On this compact, almost-horizontal chart, intersection with \(x=0\) and direction are uniformly bi-Lipschitz coordinates for the lines. Consequently, all pairs \((a,b)\) whose tubes are contained in a fixed \(r\)-tube lie in a parameter rectangle with side lengths \(O(r)\). The lattice count is \[\#\{T:T\subset T_r\}
 \leq C\left(1+\frac r{M\delta}\right)^2
 \leq C_0 r^2 N
 \leq C_1 |T_r|N.\] For \(r\) bounded below the same estimate follows by increasing \(C_1\). For \(r<\delta\), containment is impossible because each \(\delta\)-tube contains a ball of radius \(\delta\). The constants are independent of \(\delta\). For any proposed \(\eta>0\), choose \(\delta\) so small that \(C_1\leq\delta^{-\eta}\). The full shading is uniformly 1-dense. Thus every substantive hypothesis of the unrestricted Proposition 7.1 holds.\par
Projection. Choose \(f(z)=z\), so \(f'=1\) and \(f''=0\). All tubes lie in \(|z|\leq\delta\), with bounded \(x,y\). Therefore \[\bigcup_T\pi_f(Y(T))\subset[-1,1]\times[-\delta,\delta],
 \qquad \left|\bigcup_T\pi_f(Y(T))\right|\leq4\delta .\] For \(\varepsilon=1/2\) and sufficiently small \(\delta\), this is less than \(\delta^\varepsilon\), contrary to Proposition 7.1.\par
The same construction refutes Lemma 7.12 as unrestricted. For any nonempty projected subshading \(E\) inside this strip and \(r\gg\delta\), Fubini gives \[|E|\leq 2\delta\,|\operatorname{proj}_1E|,
 \qquad |N_rE|\geq r\,|\operatorname{proj}_1E|,
 \qquad |E|/|N_rE|\leq2\delta/r.\] Here the second bound follows by including the vertical interval \([-r/2,r/2]\) above each first coordinate, since \(\delta<r/2\). If \(r>\delta^{1-\varepsilon^2}\) and \(0<\varepsilon<1\), then \(2(\delta/r)^{1-\varepsilon}\to0\), contradicting (7.2). Positive shading density ensures a positive-area projection.\par
\section*{F1 continued: counterevidence, normalization and repair limits}
Counterevidence and precise scope. The horizontal example does not satisfy the Convex Wolff Axioms with small error: a slab of volume \(O(\delta)\) contains every tube. It is therefore not a counterexample to Proposition 6.8 or Theorem 1.5. The earlier sticky-Kakeya paper explicitly uses a restricted line chart with final direction component at least \(1/2\). If that full convention is intended to be inherited here, it must be stated and propagated; the counterexample is outside that revised domain.\par
A general rotation cannot justify the p. 43 reduction for a fixed \(\pi_f(x,y,z)=(x+f(z)y,z)\). Rotating a horizontal family to a vertical one changes its relation to the distinguished \(z\) coordinate and changes the projection problem. Tube Wolff non-concentration alone does not supply the missing transversality.\par
A second domain check at p. 44. The sentence invoking [16, Lemma 7.3] asserts that \[M_f=\{a+bf(t)+dtf(t):a,b,d\in[-1,1]\}\] is cinematic. The cited input has the additional normalization \(f(0)=0\). The derivative conditions printed in Proposition 7.1 alone do not imply this cinematic assertion. Take \[f(t)=450+\frac32t+\frac{t^2}{200},\quad -1\leq t\leq1.\] It satisfies \(1\leq f'(t)\leq2\) and \(|f''|=1/100\). But the two members \(0\) and \[g(t)=1-\frac{f(t)}{450}+\frac{tf(t)}{135000}
     =\frac{t^3}{27000000}\] belong to \(M_f\), are distinct, and have identical value, first derivative and second derivative at zero. Hence \[\inf_t\bigl(|g(t)|+|g'(t)|+|g''(t)|\bigr)=0
 \quad\hbox{while}\quad \|g\|_{C^2}>0.\] This violates the required cinematic separation inequality for every finite constant. It disproves this unrestricted input claim, not by itself the projected-area conclusion.\par
The exact jet determinant for the basis \(1,f,tf\) is \[\det\begin{pmatrix}
 1&f&tf\\0&f'&f+tf'\\0&f''&2f'+tf''
 \end{pmatrix}=2(f')^2-ff''.\] With \(f(0)=0\) and the stated derivative bounds on \([-1,1]\), \(|f|\leq2\) and this determinant is at least \(1.98\). This confirms why the normalized input avoids the displayed obstruction.\par
Repair boundary. State a projection lemma for a fixed transverse chart, for example \(|\operatorname{dir}(T)\cdot e_3|\geq c>0\), bounded line parameters, and a normalized, bounded slope function. Track any change of coordinates in both the tube family and the projection. Then establish that the extremal family delivered by Section 6 has these properties in those same coordinates. The prior restricted theorem is relevant counterevidence and a plausible repair route; it is not a proof of the larger statement as printed. A complete quantitative reconstruction of that interface, including all transferred refinements, was not completed here.\par
Impact. These domain errors occur in the advertised generalization used to close the generalized sticky-Kakeya argument. They are Major under the review standard even if the intended restricted route can be recovered. No assertion is made that the principal dimension conclusions are false.\par
\section*{F2. The measure-error exponent does not yield a contradiction}
Location. Proof of Theorem 1.4, p. 46: the choice \(\varepsilon=(3-\alpha)/4\), \(\eta_1=(3-\varepsilon)\eta/6\), and the last display. The identical choices occur in the acquired TeX.\par
Exact calculation. Write \(d=3-\varepsilon-\alpha=3\varepsilon>0\) and temporarily write \(\kappa\) for the measure-error exponent. The covering-number and two-sided measure estimates give, after the usual fixed enlargements of balls, \[(r/\rho)^{3-\varepsilon}\rho^{\alpha+\kappa}
 \lesssim r^{\alpha-\kappa},
 \qquad
 (r/\rho)^d\lesssim(r\rho)^{-\kappa}.\] Since \(\delta\leq\rho\leq\delta^\eta r\) and \(r\geq\delta\), the available comparison is \[\delta^{-\eta d}\lesssim\delta^{-2\kappa}.\] A contradiction requires \(\eta d>2\kappa\). For the printed choice this margin is \[\eta(3-\varepsilon-\alpha)-\frac{\eta(3-\varepsilon)}3
 =\eta\left(\frac{10\varepsilon}{3}-1\right).\] It is negative whenever \(\varepsilon<3/10\), which includes the dimension range near 3 that the argument must exclude.\par
For a concrete numerical check take \[\alpha=\frac52,\quad \varepsilon=\frac18,\quad
 d=\frac38,\quad \kappa_{\rm printed}=\frac{23\eta}{48}.\] The two exponents are \(\eta d=3\eta/8\) and \(2\kappa_{\rm printed}=23\eta/24\). The claimed contradiction is unavailable. For example \(\rho=\delta\), \(r=\delta^{1-\eta}\), with \(\eta=1/10\), meets the scale conditions and makes the preceding mass comparison compatible. This is a countercheck to the inference, not a proposed Kakeya set.\par
The first equality in the last display is also false: \(d=3-\varepsilon-\alpha\) is generally not \((3-\varepsilon)/2\). The power of \(r/\rho\) arising from rearrangement is \(+d\), not the printed \(-\eta d\). Correcting those display errors alone does not fix the overly large value of \(\eta_1\).\par
Complete local repair, conditional on Theorem 1.5. Keep \(\varepsilon=(3-\alpha)/4\), obtain \(\eta\) from Theorem 1.5, and choose \[\kappa=\frac{\eta(3-\varepsilon-\alpha)}4=\frac{3\eta\varepsilon}{4}>0.\] Definition 1.3 permits this arbitrarily small error exponent. Choose its measure first, then a positive-measure direction subset with a uniform positive line-intersection length. The associated tube shading has fixed positive density, and its Convex Wolff constant is fixed; both meet Theorem 1.5 once \(\delta\) is small enough.\par
To apply the measure bounds at legitimate centers, choose a separated subset of \(\operatorname{supp}\mu\) within \(O(\rho)\) of the points supplied by the \(\rho\)-neighborhood estimate. A fixed-factor thinning gives \(\gtrsim(r/\rho)^{3-\varepsilon}\) separated support points in an enlarged \(O(r)\) ball. Their small balls are disjoint, and the two-sided measure bounds yield exactly the comparison above, with fixed constants. Consequently \[\delta^{-\eta d}\lesssim\delta^{-2\kappa}
 =\delta^{-\eta d/2},\] which is impossible as \(\delta\downarrow0\). Thus the local deduction of Theorem 1.4 is repaired without strengthening its hypotheses.\par
Assessment. Major for the invalid final inference in a stated main theorem; repair complete at this local interface. This does not certify the independent proof of Theorem 1.5 or resolve F1.\par
\section*{Coverage and positive verification}
PDF pages | Scope actually read and checked\par
1–5 | Dimension notions, stable equality, quantifier order, measurable shadings, scale separation and the main proof architecture.\par
6–9 | Regularization, convex versus tube non-concentration, uniform partitioning covers, rescaling and the distinction between all-scale and self-similar conditions.\par
10–12 | Supremum/limsup definitions of \(\omega\) and \(\alpha\), monotonicity and existence of near-extremal pairs.\par
13–24 | Robust broadness, localization, transverse plank intersections, \(L^2\) overlap sums, enlargement, slab bound, and transfer back to tubes.\par
25–31 | Multiscale regularization, four-way alternative, elimination of richer counterexamples, anisotropic density transfer, and the reduction to Theorem 5.2.\par
32–38 | Imported sticky-Kakeya structure, balanced covers, local/global grains, slope regularity and the replacement of direction separation by slab non-concentration.\par
39–45 | Twisted projections, Frostman/Katz–Tao conversions, uniform refinement, scale merging, translated copies, cinematic input and iteration. F1.\par
45–48 | Assouad deduction, stable-equality measure argument (F2), two-scale corollary, and reference correspondence.\par
Extremal-parameter check. At a fixed point there are at most \(O(\delta^{-2})\) essentially distinct \(\delta\)-tubes in a bounded region: angular cells have area \(\asymp\delta^2\), and only boundedly many positions in a cell can pass through the point. If \(N\geq\delta^{-u+\varepsilon}\) and the shading is \(\delta^\varepsilon\)-dense, \[\left|\bigcup_TY(T)\right|\gtrsim
 \frac{\delta^\varepsilon N\delta^2}{\delta^{-2}}
 \gtrsim\delta^{4-u+2\varepsilon}.\] Taking \(r=1\), \(\rho=\delta\) in the defining local estimate supports \(\omega(u)\leq4-u\) for \(2\leq u\leq4\), and \(\omega(4)=0\). In particular, if \(\omega>0\) then \(\alpha\leq4-\omega<4\). A uniform gap in Proposition 5.1 can be selected using \(\beta+\tau\geq\alpha+\tau/2\), after choosing the loss smaller than \(\tau/4\), rather than making constants depend on a varying \(\beta\).\par
Overlap and rescaling checks. For a fixed \(s\times1\times1\) slab the dyadic overlap sum is \(\lesssim Ks^2N\log(1/s)\). Summing over all \(N\) slabs and applying Cauchy–Schwarz to total shading mass \(K^{-1}sN\) gives \(|E|\gtrsim K^{-3}/\log(1/s)\), as required in Lemma 4.4. The repeated \(N\) factor omitted in the typeset intermediate identity is not treated as an independent scientific defect.\par
The inverse tube rescaling has Jacobian \(\rho^2\). A ball of radius \(r_2\) has preimage dimensions \(\rho r_2,\rho r_2,r_2\). The regular-shading argument loses a positive density factor \(\delta^\eta\), in agreement with (5.8), rather than gaining \(\delta^{-\eta}\) as printed in (5.6). These dimensional corrections preserve the intended exponent-slack mechanism; a new global proof is not claimed.\par
\section*{Limits, interpretation and evidence}
Resolved or excluded observations. Reversed scale ratios in Definition 7.2 and the corollary’s final display, the sign of the Wolff error in Proposition 4.1, and missing scale factors in several intermediate displays were not counted separately. Their intended conventions are determined by neighboring definitions and the subsequent arguments. The iteration bound on p. 45 also needs a harmless quantitative adjustment: the recurrence \(a_{k+1}\leq(1-\varepsilon^2)a_k\) reaches \(a_k\leq\varepsilon^2\) after \[N_\varepsilon=\left\lceil
 \frac{\log(\varepsilon^2)}{\log(1-\varepsilon^2)}
 \right\rceil\] steps, rather than necessarily \(\varepsilon^{-2}\) steps. This is still a finite number depending only on \(\varepsilon\), so the initial loss can be chosen small enough. For \(\varepsilon=1/10\) the displayed formula gives 459 steps, whereas 100 is not a valid worst-case bound.\par
External-input inspection. Selected passages of Wang–Zahl, Sticky Kakeya sets and the sticky Kakeya conjecture, https://arxiv.org/html/2210.09581v2, were read: the restricted line convention, the projection lemma for AD sets, the grain-to-slab argument, and the cinematic definition, estimate and normalized-family proof. That accessible v2 is dated 7 October 2025, later than the bound manuscript. It is supporting comparison evidence, not a verified substitute for the exact earlier version cited here. Attempts to retrieve the earlier v1 failed.\par
Not fully verified. The full external grain-structure and radial-projection proofs, all scale-merging inputs, and the cinematic maximal estimate were not independently re-proved. The exact constants and nesting construction in Lemma 2.13, and the weighted selection/truncation details in the proof of Corollary 1.7, require a fuller reconstruction. Local geometric suspicions about these steps are not recorded as additional confirmed defects. The complete corrected chain from all-scale Convex Wolff Axioms to Theorem 1.5 has not been certified by this initial review.\par
Aspect | Initial-review position\par
Judgment | Major: false unrestricted supporting statements and an invalid main-theorem exponent comparison.\par
Repair | Restrict and normalize the projection input, verify its application in the same coordinates; replace the measure-error exponent as proved above.\par
Adoption | Do not use Proposition 7.1 or Lemma 7.12 in the unrestricted form. The main dimension claims are not refuted by these examples.\par
Remaining | Independent assessment of the corrected domain transfer and the explicitly unverified proof interfaces. Initial delivery is not scientific acceptance.\par
Binding and actual evidence. The primary PDF has 538,831 bytes and SHA-256\par
`0e6ba3cfaa06f55ac4fea715c21a5d51e4c1dc4838b31703ca79a03bd6b42042`.\par
All five manifest files were rehashed at intake. The 2,040-line, 189,171-byte TeX was parsed statically; parsing was not counted as scientific reading. Twenty-four original pages were actually viewed as images in addition to the complete text reading. Every final report page is separately required to pass actual visual inspection. Fifteen reviewer-written checks passed with Python 3.13.5 and SymPy 1.13.3. An initial SymPy Boolean-summation error and its corrected rerun are retained.\par
This report source was compiled using the existing authorized Tectonic 0.17.0 in untrusted mode. Actual logs, exact source/compiler/PDF hashes, dependency rules, a conservative cache inventory and page images accompany it. The cache inventory may include unused resources; it is not an exact consumed-file trace. No author code or author TeX was executed.\par
\end{document}
